\chapter{Ответы на контрольные вопросы}

\section*{Вопрос 1. Какая библиотека Python позволит составить графики любого типа?}
\addcontentsline{toc}{section}{Вопрос 1. Библиотеки Python для построения графиков}

\textbf{Ответ:}
Основным универсальным низкоуровневым инструментом для визуализации данных в экосистеме Python является библиотека \textbf{Matplotlib} (в частности, ее интерфейс \texttt{matplotlib.pyplot}). Она предоставляет полный контроль над каждым геометрическим элементом рисунка: координатными осями, сетками, логарифмическими шкалами, цветовыми палитрами, текстовыми аннотациями и векторными примитивами. Matplotlib позволяет строить графики любых типов:
\begin{itemize}
    \item линейные графики динамики функций потерь и метрик обучения;
    \item диаграммы рассеяния (scatter plots);
    \item двумерные и трехмерные гистограммы распределений;
    \item тепловые карты и матрицы ошибок (confusion matrices);
    \item трехмерные изоповерхности и контурные графики ландшафтов оптимизации.
\end{itemize}
Поверх Matplotlib функционирует библиотека \textbf{Seaborn}, предоставляющая высокоуровневые статистические темы оформления и автоматическое построение доверительных интервалов.

\section*{Вопрос 2. Какой метод позволяет изменить размерность входных данных?}
\addcontentsline{toc}{section}{Вопрос 2. Методы преобразования размерностей тензоров}

\textbf{Ответ:}
Преобразование размерностей в задачах машинного и глубокого обучения разделяется на два принципиальных уровня:
\begin{enumerate}
    \item \textbf{Перестройка формы тензора без изменения объема данных:} метод \texttt{.reshape()} в библиотеках NumPy и TensorFlow (например, преобразование матрицы в одномерный вектор: \texttt{tensor.reshape(-1, 784)}). При этом физическое расположение данных в оперативной памяти сохраняется, а изменяются лишь шаги индексации (strides).
    \item \textbf{Геометрическое масштабирование изображений:} методы интерполяции сетки пикселей. В TensorFlow используется функция \texttt{tf.image.resize(images, [height, width])}, поддерживающая билинейную, бикубическую интерполяцию и метод ближайшего соседа. В библиотеке PIL аналогичная операция реализуется методом \texttt{Image.resize((width, height))}.
    \item \textbf{Преобразования размерности внутри архитектуры нейросети:}
    \begin{itemize}
        \item слой \texttt{Flatten()} вытягивает многомерные карты признаков $(B, H, W, C)$ в двумерный тензор $(B, H \cdot W \cdot C)$ для подачи в полносвязные слои;
        \item слои пространственного пулинга (\texttt{MaxPooling2D}, \texttt{AveragePooling2D}) уменьшают высоту и ширину карт признаков с заданным шагом;
        \item слой \texttt{GlobalAveragePooling2D()} усредняет каждую карту признаков по всему пространству, сокращая размерность до $(B, C)$.
    \end{itemize}
\end{enumerate}

\section*{Вопрос 3. Чем отличается свёрточная модель от полносвязной модели нейронной сети?}
\addcontentsline{toc}{section}{Вопрос 3. Отличия сверточных и полносвязных моделей}

\textbf{Ответ:}
Фундаментальные различия заключаются в организации связей и предположениях о структуре данных:
\begin{enumerate}
    \item \textbf{Локальность рецептивных полей:} в полносвязной сети (Dense / MLP) каждый нейрон связан абсолютно со всеми входными признаками предыдущего слоя. В сверточной сети нейрон связан исключительно с малой пространственной окрестностью (размер ядра $3\times3$ или $5\times5$), что отражает пространственную коррелированность соседних пикселей.
    \item \textbf{Разделение весов (Weight Sharing):} в сверточном слое одно и то же ядро весов сканирует все изображение скользящим окном. Это на порядки сокращает число обучаемых весов: для слоя Conv2D с 64 фильтрами $3\times3$ на входе RGB требуется всего $3 \times 3 \times 3 \times 64 = 1728$ весов независимо от разрешения исходной фотографии. В полносвязном слое при входе $96\times96\times3$ потребовалось бы более 1.7 миллиона связей на каждый скрытый нейрон.
    \item \textbf{Эквивариантность к параллельному переносу:} сдвиг объекта на изображении приводит к эквивалентному сдвигу в карте активаций сверточного слоя. Полносвязная сеть вынуждена заново обучаться распознавать объект в каждом возможном положении.
    \item \textbf{Сохранение топологической структуры:} свертка естественным образом обрабатывает 2D/3D тензоры каналов, не разрушая пространственные паттерны.
\end{enumerate}

\section*{Вопрос 4. Каких видов бывают нейронные сети?}
\addcontentsline{toc}{section}{Вопрос 4. Классификация архитектур нейронных сетей}

\textbf{Ответ:}
В соответствии с лекционным курсом и современной систематикой глубокого обучения выделяют следующие базовые классы:
\begin{itemize}
    \item \textbf{Сети прямого распространения (Feedforward / MLP):} простые персептроны и многослойные полносвязные сети, где передача сигнала осуществляется строго от входа к выходу без обратных петель. Применяются для табличных данных и регрессионного анализа.
    \item \textbf{Свёрточные нейронные сети (CNN):} архитектуры со слоями локальных сверток и пулинга, специализированные для обработки сигналов с пространственной координатной сеткой (изображения, видео, спектрограммы).
    \item \textbf{Рекуррентные нейронные сети (RNN, LSTM, GRU):} сети с обратными циклическими связями («памятью»), предназначенные для обработки последовательностей переменной длины и временных рядов (текст, аудио, сигнатурный анализ).
    \item \textbf{Трансформеры (Transformers):} бескуррентные модели на базе механизма внимания (Self-Attention), способные параллельно обрабатывать контекстные зависимости произвольной дальности. Являются стандартом в обработке естественного языка (LLM) и активно применяются в компьютерном зрении (Vision Transformers, ViT).
\end{itemize}
Правильный вариант ответа в лекционном тесте кафедры: \textbf{в) Рекуррентные, сверточные, полносвязные и трансформеры}.
